\documentclass[11pt]{article}

% --- Fonts & encoding ---
\usepackage{lmodern}
\usepackage[T1]{fontenc}

% --- Page layout ---
\usepackage{geometry}
\geometry{margin=1.15in,headheight=14pt}

% --- Mathematics ---
\usepackage{amsmath,amssymb,amsthm}

% --- Typography ---
\usepackage{microtype}
\usepackage{enumitem}

% --- Tables & figures ---
\usepackage{booktabs}
\usepackage{array}

% --- Colours & hyperlinks ---
\usepackage{xcolor}
\usepackage[colorlinks,linkcolor=blue!60!black,citecolor=green!50!black,urlcolor=blue!70!black]{hyperref}

% --- Headers ---
\usepackage{fancyhdr}
\pagestyle{fancy}
\fancyhf{}
\fancyhead[L]{\small\itshape Compact Sigma Protocols for CMPT}
\fancyhead[R]{\small\thepage}
\renewcommand{\headrulewidth}{0.4pt}
\fancypagestyle{plain}{\fancyhf{}\fancyfoot[C]{\small\thepage}\renewcommand{\headrulewidth}{0pt}}

% --- Bibliography ---
\usepackage[numbers,sort&compress]{natbib}

% --- Cross-references ---
\usepackage[capitalise,noabbrev]{cleveref}
\crefname{remark}{Remark}{Remarks}
\crefname{definition}{Definition}{Definitions}
\crefname{protocol}{Protocol}{Protocols}
\crefname{equation}{Equation}{Equations}
\creflabelformat{equation}{#2(#1)#3}

% ===== Theorem environments =====
\newtheorem{theorem}{Theorem}[section]
\newtheorem{lemma}[theorem]{Lemma}
\newtheorem{corollary}[theorem]{Corollary}
\newtheorem{proposition}[theorem]{Proposition}
\theoremstyle{definition}
\newtheorem{definition}[theorem]{Definition}
\newtheorem{remark}[theorem]{Remark}
\newtheorem{protocol}[theorem]{Protocol}

% ===== Notation =====
\newcommand{\G}{\mathbb{G}}
\newcommand{\Zq}{\mathbb{Z}_q}
\newcommand{\pk}{\mathit{pk}}
\newcommand{\sk}{\mathit{sk}}
\newcommand{\Enc}{\mathsf{Enc}}
\newcommand{\Dec}{\mathsf{Dec}}
\newcommand{\Setup}{\mathsf{Setup}}
\newcommand{\KeyGen}{\mathsf{KeyGen}}
\newcommand{\EncZero}{\mathsf{EncZero}}
\newcommand{\pp}{\mathit{pp}}
\newcommand{\negl}{\mathsf{negl}}
\newcommand{\Adv}{\mathsf{Adv}}
\newcommand{\TransCtx}{\textsf{TransactionContextID}}

% ===== Title =====
\title{%
  Compact Sigma Protocols for Standard EC-ElGamal\\
  in Confidential Multi-Purpose Tokens on XRPL}
\author{Tess Dore \and Murat Cenk\\[2pt]\normalsize RippleX Research}
\date{\today}

\begin{document}
\maketitle

\begin{abstract}
  The Confidential Multi-Purpose Token (CMPT) protocol on the
  XRP Ledger~\cite{CenkMalhotraAkinyele2026} uses standard EC-ElGamal
  encryption with three separate sigma protocols---same-plaintext
  equality, amount linkage (Variant~A), and balance linkage
  (Variant~B)---contributing approximately $586$~bytes to each
  \textsf{ConfidentialMPTSend} transaction.
  %
  We show that AND-composing these three protocols under a shared
  Fiat--Shamir challenge~\cite{CramerDamgardSchoenmakers1994} and
  transmitting the proof in \emph{compact form} (challenge plus
  response scalars, rather than first-round commitments) reduces the
  sigma-proof overhead to \textbf{192~bytes}: one $32$-byte challenge
  and five $32$-byte responses.
  The verifier reconstructs all commitments algebraically and
  recomputes the hash.
  %
  Combined with the shared-$C_1$ ciphertext layout, this reduces the
  total \textsf{ConfidentialMPTSend} transaction to
  $\mathbf{1177}$~bytes for $n = 4$ recipients (the production
  configuration: sender, receiver, issuer, auditor), or $1144$~bytes
  for $n = 3$---a $27$--$29\%$ reduction from the $1604$-byte baseline,
  with no change to security assumptions.
  %
  We provide full security proofs (completeness, $2$-special soundness,
  HVZK) and implementation guidance (deterministic nonces,
  Fiat--Shamir domain separation, batch verification tradeoffs).
  %
  We also present a constant-size ($64$-byte) proof of
  shared-randomness plaintext equality across $n$ ciphertexts, due to
  Cenk~\cite{Cenk2026unification}, that may be useful for
  multi-auditor extensions.
  %
  In \cref{app:twisted}, we evaluate the twisted ElGamal variant and
  show that, under compact form, its advantage over standard EC-ElGamal
  shrinks to an irreducible $64$~bytes---insufficient to justify the
  specification, implementation, and audit costs.
\end{abstract}

\tableofcontents
\medskip

%--------------------------------------------------------------------
\section{Introduction}
\label{sec:intro}
%--------------------------------------------------------------------

The CMPT specification~\cite{CenkMalhotraAkinyele2026} defines a
\textsf{ConfidentialMPTSend} transaction whose zero-knowledge proof
bundle consists of three components:
\begin{enumerate}[label=(\roman*)]
  \item An aggregated Bulletproof~\cite{Bunz2018bulletproofs} range
        proof on the transfer amount and remainder balance
        ($754$~bytes).
  \item A same-plaintext equality proof $\Pi_{\mathrm{equal}}$ that all
        ciphertexts encrypt the same amount ($229$~bytes).
  \item Two linkage proofs: Variant~A (amount commitment to ciphertext,
        $195$~bytes) and Variant~B (balance commitment to balance
        ciphertext, $195$~bytes).
\end{enumerate}
Components (ii) and (iii) are standard Sigma protocols made
non-interactive via Fiat--Shamir.
They share overlapping witnesses---the transfer amount~$m$, shared
randomness~$r$, and sender secret key~$\sk_A$ appear in multiple
relations---but the specification treats them as independent proofs
with separate challenges and separate first-round commitments.

We observe two compression opportunities that, combined, reduce the
sigma-proof overhead from $586$ to $192$~bytes:

\begin{description}[style=nextline,leftmargin=1.5em]
  \item[AND-composition.]
    All three sigma protocols prove statements about the same
    underlying witness.
    Merging them under a single Fiat--Shamir
    challenge~\cite{CramerDamgardSchoenmakers1994} eliminates
    duplicate nonces and responses for shared witnesses, reducing
    the merged proof to $n + 5$ group-element commitments
    and $5$ scalar responses.

  \item[Compact (challenge) form.]
    In a Fiat--Shamir Sigma protocol, the verifier can
    \emph{reconstruct} the first-round commitments from the
    challenge and responses via the verification equations, then
    recompute the hash to check consistency.
    This is standard folklore---Schnorr
    signatures~\cite{Schnorr1991} transmit $(e, z)$ rather than
    $(R, z)$.
    Applying this to the AND-composed protocol eliminates all
    $n + 5$ group elements from the proof, leaving only $6$
    scalars: one challenge and five responses.
\end{description}

The resulting compact proof is $6 \times 32 = 192$~bytes, regardless
of the number of recipients~$n$.
Combined with the shared-$C_1$ ciphertext deduplication,
the total \textsf{ConfidentialMPTSend} size decreases from $1604$ to
$1177$~bytes for $n = 4$ (sender, receiver, issuer, auditor---the
production configuration) or $1144$~bytes for $n = 3$.

\medskip
\noindent\textbf{Scope and structure.}
\Cref{sec:relation} derives the combined proof statement from the
CMPT specification.
\Cref{sec:protocol} presents the AND-composed sigma protocol in both
interactive and compact forms.
\Cref{sec:security} provides full security proofs.
\Cref{sec:equality-o1} presents a constant-size shared-randomness
equality proof due to Cenk~\cite{Cenk2026unification}.
\Cref{sec:implementation} discusses implementation considerations.
\Cref{sec:sizes} gives the transaction size analysis.
\Cref{app:twisted} evaluates the twisted ElGamal variant and shows
that compact form reduces its advantage to $64$~bytes.

%--------------------------------------------------------------------
\section{Combined Proof Statement}
\label{sec:relation}
%--------------------------------------------------------------------

We derive the combined relation from the CMPT specification,
following~\cite{Cenk2026unification}.

\subsection{Setup and Notation}

Let $\G$ be the secp256k1 group of prime order~$q$ with base
point~$G$.
Let $H \in \G$ be an independent generator (NUMS point) for Pedersen
commitments, with $\log_G H$ unknown.
Let $m \in \Zq$ denote the transfer amount and
$r \in \Zq$ the shared encryption randomness.
For $n$ public keys $P_1, \ldots, P_n \in \G$---where $P_A$ denotes
the sender's public key---the transfer amount is represented by a
shared-randomness EC-ElGamal ciphertext family.
In production, $n = 4$ (sender, receiver, issuer, auditor);
we give sizes for both $n = 3$ and $n = 4$ throughout.
\[
  C_1 = r \cdot G, \qquad
  C_{2,i} = m \cdot G + r \cdot P_i, \quad i = 1, \ldots, n.
\]
The shared first component~$C_1$ is transmitted once on-chain.

\subsection{Pedersen Commitments}

The same transfer amount is committed using a Pedersen commitment
that intentionally reuses the ciphertext randomness~$r$:
\[
  \mathit{PC}_m = m \cdot G + r \cdot H.
\]
This eliminates the need for a separate Variant~A linkage witness:
the sigma proof naturally handles the linkage because $m$ and $r$
appear in both the ciphertext and commitment relations.

The sender's balance before transfer is $b \in \Zq$, committed with
fresh randomness $\rho \in \Zq$:
\[
  \mathit{PC}_b = b \cdot G + \rho \cdot H.
\]

\subsection{Balance Ciphertext Linkage}

Let $(B_1, B_2)$ denote the sender's existing on-ledger balance
ciphertext under the sender's public key~$P_A$.
For some historical randomness $r_b \in \Zq$:
\[
  B_1 = r_b \cdot G, \qquad
  B_2 = b \cdot G + r_b \cdot P_A.
\]
Since $(B_1, B_2)$ is already on-ledger, the prover may not know the
original $r_b$.
However, the sender knows $\sk_A \in \Zq$ with $P_A = \sk_A \cdot G$,
so the balance ciphertext satisfies:
\[
  B_2 - b \cdot G = \sk_A \cdot B_1.
\]
This allows the sigma proof to link $\mathit{PC}_b$ to the balance
ciphertext using $\sk_A$ rather than~$r_b$.

\subsection{Combined Relation}

The prover's witness is
\[
  w = (m,\; r,\; b,\; \rho,\; \sk_A).
\]
The combined relation for \textsf{ConfidentialMPTSend} is
\begin{equation}
  \label{eq:combined-relation}
  \mathcal{R}_{\mathrm{send}} = \left\{
  \bigl(
    \mathbf{x},\; w
  \bigr)
  \;\middle|\;
  \begin{aligned}
    C_1 &= r \cdot G, \\
    C_{2,i} &= m \cdot G + r \cdot P_i
      \quad (i = 1, \ldots, n), \\
    \mathit{PC}_m &= m \cdot G + r \cdot H, \\
    \mathit{PC}_b &= b \cdot G + \rho \cdot H, \\
    P_A &= \sk_A \cdot G, \\
    B_2 - b \cdot G &= \sk_A \cdot B_1
  \end{aligned}
  \right\},
\end{equation}
where $\mathbf{x} = (C_1, \{C_{2,i}\}_{i=1}^n, \mathit{PC}_m,
\mathit{PC}_b, B_1, B_2, P_1, \ldots, P_n, P_A, H)$.

This relation captures three consistency requirements simultaneously:
(i) all amount ciphertexts encrypt the same value~$m$ using shared
randomness~$r$;
(ii) the Pedersen commitment $\mathit{PC}_m$ commits to the same
amount~$m$;
and (iii) the Pedersen commitment $\mathit{PC}_b$ commits to the same
balance~$b$ as the existing balance ciphertext.

The aggregated Bulletproof separately proves $m \in [0, 2^{64})$ and
$b - m \in [0, 2^{64})$ using commitments $\mathit{PC}_m$ and
$\mathit{PC}_b - \mathit{PC}_m = (b - m) \cdot G + (\rho - r) \cdot H$.

%--------------------------------------------------------------------
\section{AND-Composed Compact-Form Sigma Protocol}
\label{sec:protocol}
%--------------------------------------------------------------------

\subsection{Interactive Protocol}
\label{sec:interactive}

\begin{protocol}[AND-composed Sigma protocol $\Sigma_{\mathrm{send}}$]
  \label{prot:combined}
  On public input~$\mathbf{x}$ and private input
  $w = (m, r, b, \rho, \sk_A)$:

  \begin{enumerate}
    \item \textbf{Commit.}
          Prover samples nonces
          $\alpha_m, \alpha_r, \alpha_b, \alpha_\rho,
           \alpha_{\sk} \xleftarrow{\;R\;} \Zq$
          and computes $n + 5$ commitment elements:
          \begin{align}
            T_1 &= \alpha_r \cdot G,
              \label{eq:com-t1} \\
            T_{2,i} &= \alpha_m \cdot G + \alpha_r \cdot P_i
              \quad (i = 1, \ldots, n),
              \label{eq:com-t2} \\
            T_m &= \alpha_m \cdot G + \alpha_r \cdot H,
              \label{eq:com-tm} \\
            T_b &= \alpha_b \cdot G + \alpha_\rho \cdot H,
              \label{eq:com-tb} \\
            T_{\sk,1} &= \alpha_{\sk} \cdot G,
              \label{eq:com-tsk1} \\
            T_{\sk,2} &= \alpha_b \cdot G + \alpha_{\sk} \cdot B_1.
              \label{eq:com-tsk2}
          \end{align}
          Prover sends
          $(T_1, T_{2,1}, \ldots, T_{2,n}, T_m, T_b,
            T_{\sk,1}, T_{\sk,2})$.

    \item \textbf{Challenge.}
          Verifier sends $e \xleftarrow{\;R\;} \Zq$.

    \item \textbf{Response.}
          Prover computes and sends:
          \begin{align}
            z_m &= \alpha_m + e \cdot m,
              \label{eq:resp-zm} \\
            z_r &= \alpha_r + e \cdot r,
              \label{eq:resp-zr} \\
            z_b &= \alpha_b + e \cdot b,
              \label{eq:resp-zb} \\
            z_\rho &= \alpha_\rho + e \cdot \rho,
              \label{eq:resp-zrho} \\
            z_{\sk} &= \alpha_{\sk} + e \cdot \sk_A.
              \label{eq:resp-zsk}
          \end{align}

    \item \textbf{Verification.}
          Verifier accepts iff all of the following hold:
          \begin{align}
            z_r \cdot G
              &= T_1 + e \cdot C_1,
              \label{eq:ver-t1} \\
            z_m \cdot G + z_r \cdot P_i
              &= T_{2,i} + e \cdot C_{2,i}
              \quad \forall\, i,
              \label{eq:ver-t2} \\
            z_m \cdot G + z_r \cdot H
              &= T_m + e \cdot \mathit{PC}_m,
              \label{eq:ver-tm} \\
            z_b \cdot G + z_\rho \cdot H
              &= T_b + e \cdot \mathit{PC}_b,
              \label{eq:ver-tb} \\
            z_{\sk} \cdot G
              &= T_{\sk,1} + e \cdot P_A,
              \label{eq:ver-tsk1} \\
            z_b \cdot G + z_{\sk} \cdot B_1
              &= T_{\sk,2} + e \cdot B_2.
              \label{eq:ver-tsk2}
          \end{align}
  \end{enumerate}
\end{protocol}

This is the standard AND-composition~\cite{CramerDamgardSchoenmakers1994}
of the same-plaintext, amount-linkage, and balance-linkage protocols
under a shared challenge.
The interactive proof transmits $n + 5$ group elements and $5$ scalars.

\medskip
\noindent\textbf{Uncompressed proof size.}
The merged proof transmits $(n + 5)$ group elements and $5$ scalars:
\[
  (n + 5) \times 33 + 5 \times 32 \;=\;
  \begin{cases}
    424 \text{~bytes} & (n = 3), \\
    457 \text{~bytes} & (n = 4).
  \end{cases}
\]
This compares to $586$~bytes ($n = 3$) or $619$~bytes ($n = 4$) for
three separate proofs---a $28\%$ reduction from witness sharing alone.

\subsection{Compact (Challenge) Form}
\label{sec:compact}

Each verification equation can be rearranged to express the commitment
in terms of the challenge and responses:
\begin{align}
  T_1 &= z_r \cdot G - e \cdot C_1,
    \label{eq:recon-t1} \\
  T_{2,i} &= z_m \cdot G + z_r \cdot P_i - e \cdot C_{2,i}
    \quad \forall\, i,
    \label{eq:recon-t2} \\
  T_m &= z_m \cdot G + z_r \cdot H - e \cdot \mathit{PC}_m,
    \label{eq:recon-tm} \\
  T_b &= z_b \cdot G + z_\rho \cdot H - e \cdot \mathit{PC}_b,
    \label{eq:recon-tb} \\
  T_{\sk,1} &= z_{\sk} \cdot G - e \cdot P_A,
    \label{eq:recon-tsk1} \\
  T_{\sk,2} &= z_b \cdot G + z_{\sk} \cdot B_1 - e \cdot B_2.
    \label{eq:recon-tsk2}
\end{align}

In the Fiat--Shamir transform, the challenge is:
\begin{multline}
  \label{eq:fs-hash}
  e = \mathcal{H}\!\bigl(
    \texttt{"CMPT\_SEND\_SIGMA"} \;\|\;
    P_1 \| \cdots \| P_n \| P_A \;\|\;
    C_1 \| C_{2,1} \| \cdots \| C_{2,n} \;\|\; \\
    \mathit{PC}_m \| \mathit{PC}_b \| B_1 \| B_2 \;\|\;
    T_1 \| T_{2,1} \| \cdots \| T_{2,n} \| T_m
    \| T_b \| T_{\sk,1} \| T_{\sk,2} \;\|\;
    \TransCtx
  \bigr).
\end{multline}

\begin{definition}[Compact-form proof $\Pi_{\mathrm{send}}$]
  \label{def:compact}
  The non-interactive proof in compact form is
  \[
    \Pi_{\mathrm{send}} = (e,\, z_m,\, z_r,\, z_b,\, z_\rho,\, z_{\sk})
    \;\in\; \Zq^6.
  \]
  \textbf{Verification:}
  \begin{enumerate}
    \item Reconstruct all $n + 5$ commitments
          via \crefrange{eq:recon-t1}{eq:recon-tsk2}.
    \item Recompute $e' = \mathcal{H}(\cdots)$ using
          \cref{eq:fs-hash} with the reconstructed commitments.
    \item Accept iff $e' = e$.
  \end{enumerate}
\end{definition}

\noindent\textbf{Proof size.}
$6 \times 32 = \mathbf{192}$~bytes (one challenge, five responses;
all $\Zq$ scalars).
The proof size is \textbf{independent of~$n$};
verification cost remains $O(n)$ group operations.

\begin{remark}[Equivalence of commitment and compact forms]
  \label{rem:equivalence}
  The two forms define equivalent verification procedures:
  a proof is valid in one form if and only if the corresponding
  proof is valid in the other.
  In the commitment form, the verifier computes $e$ from the hash
  and checks the verification equations; in the compact form, the
  verifier reconstructs the commitments from the same equations and
  checks the hash.
  A valid commitment-form proof converts to compact form by
  dropping the commitments and prepending~$e$, and vice versa.
\end{remark}

%--------------------------------------------------------------------
\section{Security}
\label{sec:security}
%--------------------------------------------------------------------

We prove security properties of the interactive protocol
$\Sigma_{\mathrm{send}}$ (\cref{prot:combined}), then lift to the
non-interactive setting via the Fiat--Shamir transform.
Throughout, $\lambda$ denotes the security parameter,
$\G$ is the secp256k1 group of prime order
$q$ ($\approx 2^{256}$), with generators $G$ and $H$ such that
$\log_G H$ is unknown.

\subsection{Completeness}

\begin{theorem}[Completeness]
  \label{thm:completeness}
  If the prover holds a valid witness $w = (m, r, b, \rho, \sk_A)$
  for a statement $\mathbf{x} \in \mathcal{R}_{\mathrm{send}}$,
  then the verifier always accepts.
\end{theorem}

\begin{proof}
  We verify each equation.

  \medskip\noindent\textit{\Cref{eq:ver-t1}:}
  $z_r \cdot G = (\alpha_r + er) \cdot G
    = \alpha_r \cdot G + e \cdot r \cdot G
    = T_1 + e \cdot C_1.\; \checkmark$

  \medskip\noindent\textit{\Cref{eq:ver-t2}, for each $i$:}
  \begin{align*}
    z_m \cdot G + z_r \cdot P_i
    &= (\alpha_m + em) \cdot G + (\alpha_r + er) \cdot P_i \\
    &= (\alpha_m \cdot G + \alpha_r \cdot P_i)
       + e \cdot (m \cdot G + r \cdot P_i) \\
    &= T_{2,i} + e \cdot C_{2,i}.\; \checkmark
  \end{align*}

  \medskip\noindent\textit{\Cref{eq:ver-tm}:}
  $z_m \cdot G + z_r \cdot H
    = (\alpha_m \cdot G + \alpha_r \cdot H)
      + e \cdot (m \cdot G + r \cdot H)
    = T_m + e \cdot \mathit{PC}_m.\; \checkmark$

  \medskip\noindent\textit{\Cref{eq:ver-tb}:}
  $z_b \cdot G + z_\rho \cdot H
    = (\alpha_b \cdot G + \alpha_\rho \cdot H)
      + e \cdot (b \cdot G + \rho \cdot H)
    = T_b + e \cdot \mathit{PC}_b.\; \checkmark$

  \medskip\noindent\textit{\Cref{eq:ver-tsk1}:}
  $z_{\sk} \cdot G
    = (\alpha_{\sk} + e \cdot \sk_A) \cdot G
    = \alpha_{\sk} \cdot G + e \cdot \sk_A \cdot G
    = T_{\sk,1} + e \cdot P_A.\; \checkmark$

  \medskip\noindent\textit{\Cref{eq:ver-tsk2}:}
  \begin{align*}
    z_b \cdot G + z_{\sk} \cdot B_1
    &= (\alpha_b + eb) \cdot G + (\alpha_{\sk} + e \cdot \sk_A) \cdot B_1 \\
    &= (\alpha_b \cdot G + \alpha_{\sk} \cdot B_1)
       + e \cdot (b \cdot G + \sk_A \cdot B_1) \\
    &= T_{\sk,2} + e \cdot B_2.\; \checkmark
    \qedhere
  \end{align*}
\end{proof}

\subsection{Special Soundness}

\begin{theorem}[$2$-special soundness]
  \label{thm:soundness}
  Given two accepting transcripts
  \begin{align*}
    &(T_1, T_{2,i}, T_m, T_{\sk,1}, T_b, T_{\sk,2},\;
       e,\; z_r, z_m, z_{\sk}, z_\rho, z_b) \quad\text{and}\\
    &(T_1, T_{2,i}, T_m, T_{\sk,1}, T_b, T_{\sk,2},\;
       e',\; z_r', z_m', z_{\sk}', z_\rho', z_b')
  \end{align*}
  with the same first-round commitments and $e \ne e'$,
  one can efficiently extract a witness
  $(m, r, b, \rho, \sk_A)$ satisfying all relations
  of~\cref{eq:combined-relation}.
\end{theorem}

\begin{proof}
  Define $\Delta e = e - e'$, which is nonzero and hence invertible
  in~$\Zq$.  Set:
  \[
    m = \frac{z_m - z_m'}{\Delta e},\quad
    r = \frac{z_r - z_r'}{\Delta e},\quad
    b = \frac{z_b - z_b'}{\Delta e},\quad
    \rho = \frac{z_\rho - z_\rho'}{\Delta e},\quad
    \sk_A = \frac{z_{\sk} - z_{\sk}'}{\Delta e}.
  \]

  \medskip\noindent\textit{Relation 1 ($C_1 = r \cdot G$):}
  Subtracting the two instances of \cref{eq:ver-t1}:
  $(z_r - z_r') \cdot G = (e - e') \cdot C_1$,
  so $r \cdot G = C_1$. $\checkmark$

  \medskip\noindent\textit{Relation 2 ($C_{2,i} = m \cdot G + r \cdot P_i$ for all $i$):}
  Subtracting the two instances of \cref{eq:ver-t2}:
  $(z_m - z_m') \cdot G + (z_r - z_r') \cdot P_i = (e - e') \cdot C_{2,i}$,
  so $m \cdot G + r \cdot P_i = C_{2,i}$. $\checkmark$

  \medskip\noindent\textit{Relation 3 ($\mathit{PC}_m = m \cdot G + r \cdot H$):}
  From \cref{eq:ver-tm}:
  $(z_m - z_m') \cdot G + (z_r - z_r') \cdot H = (e - e') \cdot \mathit{PC}_m$,
  so $m \cdot G + r \cdot H = \mathit{PC}_m$. $\checkmark$

  \medskip\noindent\textit{Relation 4 ($\mathit{PC}_b = b \cdot G + \rho \cdot H$):}
  From \cref{eq:ver-tb}:
  $(z_b - z_b') \cdot G + (z_\rho - z_\rho') \cdot H = (e - e') \cdot \mathit{PC}_b$,
  so $b \cdot G + \rho \cdot H = \mathit{PC}_b$. $\checkmark$

  \medskip\noindent\textit{Relation 5 ($P_A = \sk_A \cdot G$):}
  From \cref{eq:ver-tsk1}:
  $(z_{\sk} - z_{\sk}') \cdot G = (e - e') \cdot P_A$,
  so $\sk_A \cdot G = P_A$. $\checkmark$

  \medskip\noindent\textit{Relation 6 ($B_2 - b \cdot G = \sk_A \cdot B_1$):}
  From \cref{eq:ver-tsk2}:
  $(z_b - z_b') \cdot G + (z_{\sk} - z_{\sk}') \cdot B_1
   = (e - e') \cdot B_2$,
  so $b \cdot G + \sk_A \cdot B_1 = B_2$. $\checkmark$

  All six relations are satisfied by the extracted witness.
\end{proof}

\subsection{Honest-Verifier Zero-Knowledge}

\begin{theorem}[HVZK]
  \label{thm:hvzk}
  $\Sigma_{\mathrm{send}}$ is special honest-verifier zero-knowledge
  for $\mathcal{R}_{\mathrm{send}}$.
\end{theorem}

\begin{proof}
  We construct a simulator $\mathcal{S}$ that, given any challenge
  $e \in \Zq$ and a valid statement (but no witness), outputs a
  transcript indistinguishable from a real one.

  \medskip\noindent\textbf{Simulator $\mathcal{S}(e)$:}
  \begin{enumerate}
    \item Sample $z_m, z_r, z_b, z_\rho, z_{\sk}
          \xleftarrow{\;R\;} \Zq$.
    \item Set commitments via the reconstruction equations
          \crefrange{eq:recon-t1}{eq:recon-tsk2}.
    \item Output the full transcript
          $(T_1, T_{2,i}, T_m, T_{\sk,1}, T_b, T_{\sk,2},\;
            e,\; z_r, z_m, z_{\sk}, z_\rho, z_b)$.
  \end{enumerate}

  \medskip\noindent\textbf{Verification.}
  By construction, the transcript satisfies all verification equations
  \crefrange{eq:ver-t1}{eq:ver-tsk2}.

  \medskip\noindent\textbf{Distribution.}
  In a real transcript, the nonces
  $\alpha_m, \alpha_r, \alpha_b, \alpha_\rho, \alpha_{\sk}$
  are uniform in~$\Zq$, and the responses are affine functions
  of the nonces.
  For any fixed $e$ and witness, the map
  $(\alpha_m, \alpha_r, \alpha_b, \alpha_\rho, \alpha_{\sk})
   \mapsto (z_m, z_r, z_b, z_\rho, z_{\sk})$
  is a bijection on~$\Zq^5$.
  Therefore the responses are uniformly distributed, and the
  commitments are deterministic functions of the statement, challenge,
  and responses.
  The simulated transcript has the same structure.
  The two distributions are identical.
\end{proof}

\subsection{Non-Interactive Zero-Knowledge in the Random Oracle Model}

\begin{theorem}[NIZK in the ROM]
  \label{thm:nizk}
  The Fiat--Shamir transform of $\Sigma_{\mathrm{send}}$
  (in either commitment or compact form) is a non-interactive
  zero-knowledge proof of knowledge for $\mathcal{R}_{\mathrm{send}}$
  in the random oracle model, under the discrete logarithm assumption
  in~$\G$.
\end{theorem}

\begin{proof}[Proof sketch]
  \textbf{Knowledge soundness (via the forking lemma).}
  Let $\mathcal{A}$ be a PPT adversary that produces an accepting
  proof with non-negligible probability~$\varepsilon$.
  Apply the general forking lemma~\cite{BellareNeven2006}:
  run $\mathcal{A}$ once to obtain an accepting transcript,
  then fork with a fresh random oracle response at the challenge
  query.
  With probability $\ge \varepsilon^2/q_H - \varepsilon/q$, both
  runs produce accepting proofs with the same commitments but distinct
  challenges.
  By \cref{thm:soundness}, the extractor recovers
  $(m, r, b, \rho, \sk_A)$.

  \medskip
  \textbf{Zero-knowledge (via RO simulation).}
  The simulator samples uniform responses, reconstructs commitments
  via \crefrange{eq:recon-t1}{eq:recon-tsk2} using a random~$e$,
  and programs the random oracle to return~$e$ on the appropriate
  input.
  Programming succeeds except with probability
  $q_H / q = \negl(\lambda)$.
\end{proof}

\begin{remark}[Long-lived secret $\sk_A$]
  \label{rem:long-lived}
  The compact form leaks no more about the long-lived key~$\sk_A$
  than the commitment form: both admit the same witness extraction
  (\cref{thm:soundness}), and the HVZK simulator (\cref{thm:hvzk})
  produces transcripts without any witness.
  In the non-interactive setting, Fiat--Shamir prevents an adversary
  from obtaining two transcripts with the same commitments but distinct
  challenges (requiring a hash collision).
  The security of $\sk_A$ reduces to the discrete logarithm
  assumption, identically to a standalone Schnorr PoK.
\end{remark}

%--------------------------------------------------------------------
\section{Constant-Size Shared-Randomness Equality Proof}
\label{sec:equality-o1}
%--------------------------------------------------------------------

The combined sigma protocol of \cref{sec:protocol} proves ciphertext
equality as part of a larger conjunction that includes the balance
linkage.
In some contexts---standalone equality attestations, or multi-auditor
scenarios where $n$ is large---it is useful to have a proof of
shared-randomness equality alone, with proof size
\emph{independent of}~$n$.
The following construction, due to
Cenk~\cite{Cenk2026unification}, achieves this.

\subsection{Algebraic Cancellation}

Let $P_1, \ldots, P_n \in \G$ be distinct public keys and let
$r \xleftarrow{R} \Zq$.
The sender encrypts a message $M = m \cdot G$ under shared
randomness, producing $n$ ciphertexts $E_i = (C, D_i)$ with
\[
  C = r \cdot G, \qquad
  D_i = M + r \cdot P_i, \quad i = 1, \ldots, n.
\]
Since $C$ is shared, all ciphertexts have a common first component.
The goal is to prove that all $D_i$ encrypt the same~$M$
without revealing~$M$ or~$r$.

Define differences relative to the first ciphertext:
\[
  \Delta P_i = P_i - P_1, \qquad
  \Delta D_i = D_i - D_1, \qquad i = 2, \ldots, n.
\]
Then $\Delta D_i = r \cdot \Delta P_i$ (the message cancels).
The problem reduces to a multi-base DLEQ: prove knowledge of~$r$ such
that $C = r \cdot G$ and
$\Delta D_i = r \cdot \Delta P_i$ for all $i \geq 2$.

\subsection{\texorpdfstring{The $(c, s)$ Variant}{The (c, s) Variant}}

Because all DLEQ relations share a \emph{single} scalar witness~$r$,
a single Schnorr response suffices.

\begin{protocol}[$\pi_{\mathrm{eq}}$: constant-size shared-$r$ equality]
  \label{prot:equality-o1}
  On public input
  $\bigl(C,\; \Delta D_2, \ldots, \Delta D_n,$
  $\Delta P_2, \ldots, \Delta P_n\bigr)$
  and private input~$r$:

  \begin{enumerate}
    \item \textbf{Commit.}
          Sample $k \xleftarrow{R} \Zq$.
          Compute $R_G = k \cdot G$ and
          $R_i = k \cdot \Delta P_i$ for $i = 2, \ldots, n$.

    \item \textbf{Challenge (Fiat--Shamir).}
          \[
            c = H\bigl(C \| \Delta D_2 \| \cdots \| \Delta D_n
            \| R_G \| R_2 \| \cdots \| R_n\bigr).
          \]

    \item \textbf{Response.}
          $s = k + c \cdot r$.

    \item \textbf{Proof.}
          $\pi_c = (c, s)$.
  \end{enumerate}
\end{protocol}

\noindent\textbf{Verification.}
The verifier reconstructs:
\[
  R_G' = s \cdot G - c \cdot C, \qquad
  R_i' = s \cdot \Delta P_i - c \cdot \Delta D_i
  \quad (i = 2, \ldots, n),
\]
and checks $c \stackrel{?}{=}
H(C \| \Delta D_2 \| \cdots \| \Delta D_n
\| R_G' \| R_2' \| \cdots \| R_n')$.

\paragraph{Proof size.}
$\pi_c = (c, s) \in \Zq^2$: \textbf{64~bytes}, independent of~$n$.
Verification cost is $\mathcal{O}(n)$ scalar multiplications, but the
proof transmitted on-chain is $\mathcal{O}(1)$.

\paragraph{Security.}
Completeness and special soundness follow from the standard Schnorr
argument applied to the multi-base DLEQ.
The single response~$s$ suffices because all $n$ relations share the
same witness~$r$; the extractor recovers
$r = (s - s') / (c - c')$ from two accepting transcripts.

\subsection{Relationship to the Combined Sigma Proof}

For the full \textsf{ConfidentialMPTSend} proof, the combined sigma
protocol of \cref{sec:protocol} already handles shared-randomness
equality as part of the larger conjunction.
The compact form gives $192$~bytes for the full proof (equality +
amount linkage + balance linkage), independent of~$n$.

The $\pi_{\mathrm{eq}}$ construction is therefore not needed for
\textsf{ConfidentialMPTSend} at small fixed~$n$.
It becomes relevant in two scenarios:
\begin{enumerate}[nosep]
  \item \textbf{Large $n$} (e.g., multi-auditor extensions where
        $n \gg 4$): while the compact sigma proof's byte count is
        already $n$-independent, its verification cost grows as
        $\mathcal{O}(n)$.
        Separating the shared-$r$ equality proof ($64$~bytes) from
        the remaining balance relations could enable more modular
        verification strategies.
  \item \textbf{Standalone equality attestations} outside of
        \textsf{ConfidentialMPTSend} (e.g., proving re-encryption
        correctness during key rotation).
\end{enumerate}

%--------------------------------------------------------------------
\section{Implementation Considerations}
\label{sec:implementation}
%--------------------------------------------------------------------

\subsection{Deterministic Nonces}
\label{sec:nonces}

The prover's nonces
$(\alpha_m, \alpha_r, \alpha_b, \alpha_\rho, \alpha_{\sk})$ must be
sampled uniformly and independently for each proof instance.
Reusing a nonce across two proofs with distinct
challenges---or across two distinct statements with the same
challenge---allows trivial witness extraction via the special
soundness equations.

Because the compact form transmits $(e, z_m, z_r, z_b, z_\rho, z_{\sk})$
rather than the commitments, a nonce-reuse attack is not detectable
from the proof alone; the resulting proofs would simply be valid proofs
of distinct statements that happen to share the underlying nonces.
An adversary observing both proofs can extract the witnesses.

\begin{remark}[Synthetic nonces]
  \label{rem:synthetic-nonces}
  We recommend \emph{deterministic nonce generation} following
  RFC~6979~\cite{RFC6979}:
  \[
    (\alpha_m, \ldots, \alpha_{\sk}) \;\leftarrow\;
    \mathsf{HKDF}\bigl(
    \sk_A \| r \| m \| b \| \rho \| \text{statement}
    \| \text{fresh entropy}
    \bigr).
  \]
  Including the full witness ensures nonces are unique per transaction.
  The optional fresh entropy provides defense-in-depth.
\end{remark}

\subsection{Fiat--Shamir Domain Separation}
\label{sec:domain-sep}

The hash input for the challenge (\cref{eq:fs-hash}) must include:
\begin{enumerate}
  \item A fixed domain-separation tag
        (\texttt{"CMPT\_SEND\_SIGMA"}).
  \item \emph{All} public-key and ciphertext components:
        the recipient keys $P_1, \ldots, P_n$, the sender key~$P_A$,
        the ciphertext bundle $C_1, C_{2,1}, \ldots, C_{2,n}$, the
        commitments $\mathit{PC}_m, \mathit{PC}_b$, and the balance
        ciphertext~$(B_1, B_2)$.
  \item All first-round commitments
        $T_1, T_{2,1}, \ldots, T_{2,n}, T_m, T_b, T_{\sk,1}, T_{\sk,2}$.
  \item The \TransCtx{} for replay protection.
\end{enumerate}
Omitting any statement element from the hash would allow
cross-protocol attacks.
The domain-separation tag prevents collisions with other
Fiat--Shamir proofs in the CMPT protocol.

\subsection{Verification Cost}
\label{sec:verification-cost}

Compact-form and commitment-form verification require the same number
of multi-scalar multiplications.
The verifier performs identical linear combinations, just in
\emph{reconstruction} order rather than \emph{check} order.
The only additional cost is a hash evaluation to compute~$e'$,
negligible compared to the group operations.

For general~$n$, verification reconstructs $n + 5$ commitments via
$n + 3$ double-exponentiations, $1$ triple-exponentiation (for
$T_m = z_m G + z_r H - e \cdot \mathit{PC}_m$),
$1$ single-exponentiation (for $T_{\sk,1}$), and $1$ hash evaluation.
At $n = 4$: $11$ multi-scalar multiplications total;
at $n = 3$: $10$.

\begin{remark}[Batch verification tradeoff]
  \label{rem:batch}
  Compact-form proofs do not support batch verification in the
  standard way (random linear combination of independent verification
  equations across transactions), because the commitment points are
  \emph{reconstructed} rather than transmitted independently.
  In commitment form, the verifier can batch the linear check
  $T + e \cdot S \stackrel{?}{=} z \cdot B$ across proofs by drawing
  random weights; in compact form, each proof's commitments depend on
  its own challenge and responses, so no cross-proof linear structure
  is available.
  In practice this tradeoff is favourable: the on-chain proof size
  reduction (from eliminating group elements) outweighs the marginal
  throughput loss, since network latency and transaction-size limits
  dominate validator cost in the XRPL ledger-close model.
\end{remark}

%--------------------------------------------------------------------
\section{Transaction Size Analysis}
\label{sec:sizes}
%--------------------------------------------------------------------

All sizes use compressed secp256k1 points ($33$~bytes) and $\Zq$
scalars ($32$~bytes).
We give results for both $n = 4$ recipients
(sender, receiver, issuer, auditor---the production configuration)
and $n = 3$ (no auditor).
The only $n$-dependent component is the ciphertext data:
each additional recipient adds one $C_{2,i}$ group element
($+33$~bytes).
The compact sigma proof is $n$-independent.

\subsection{Sigma-Proof Comparison}

\Cref{tab:sigma-comparison} compares the sigma-proof overhead across
the three proof formats.

\begin{table}[htbp]
  \centering
  \caption{%
    Sigma-proof size for \textsf{ConfidentialMPTSend}
    (standard EC-ElGamal).
    Each row shows sigma overhead, excluding ciphertext data and
    Bulletproof.  The merged and separate forms grow by one group
    element ($+33$~B) per additional recipient; the compact form is
    $n$-independent.}
  \label{tab:sigma-comparison}
  \begin{tabular}{l@{\quad}rrr}
    \toprule
    & \shortstack{Separate\\(commit.\ form)}
    & \shortstack{Merged\\(commit.\ form)}
    & \shortstack{Compact\\form} \\
    \midrule
    \multicolumn{4}{l}{\textit{$n = 4$ (production: sender, receiver, issuer, auditor)}} \\[2pt]
    \quad Group elements
    & $11 \times 33 = 363$
    & $9 \times 33 = 297$
    & $0$ \\
    \quad Scalars
    & $8 \times 32 = 256$
    & $5 \times 32 = 160$
    & $6 \times 32 = 192$ \\
    \quad\textbf{Total}
    & \textbf{619}
    & \textbf{457}
    & \textbf{192} \\[4pt]
    \multicolumn{4}{l}{\textit{$n = 3$ (no auditor)}} \\[2pt]
    \quad Group elements
    & $10 \times 33 = 330$
    & $8 \times 33 = 264$
    & $0$ \\
    \quad Scalars
    & $8 \times 32 = 256$
    & $5 \times 32 = 160$
    & $6 \times 32 = 192$ \\
    \quad\textbf{Total}
    & \textbf{586}
    & \textbf{424}
    & \textbf{192} \\
    \bottomrule
  \end{tabular}
\end{table}

\noindent
For $n = 4$, merging alone saves $619 - 457 = 162$~bytes by
eliminating duplicate commitments and responses for shared witnesses.
Compact form saves a further $457 - 192 = 265$~bytes by eliminating
all group-element commitments.
The total reduction from separate proofs to compact form is
$619 - 192 = 427$~bytes ($69\%$).
The compact-form proof size is $192$~bytes regardless of~$n$, since
only scalar responses are transmitted.

\subsection{Full Transaction Comparison}

\begin{table}[htbp]
  \centering
  \caption{%
    Full \textsf{ConfidentialMPTSend} transaction size breakdown,
    standard EC-ElGamal with compact sigma.}
  \label{tab:full-breakdown}
  \begin{tabular}{l@{\quad}rrr}
    \toprule
    Component & $n = 4$ & $n = 3$ & Notes \\
    \midrule
    Shared $C_1$
      & $33$ & $33$ & $1 \times 33$ \\
    Recipient ciphertexts $C_{2,i}$
      & $132$ & $99$ & $n \times 33$ \\
    Amount commitment $\mathit{PC}_m$
      & $33$ & $33$ & $1 \times 33$ \\
    Balance commitment $\mathit{PC}_b$
      & $33$ & $33$ & $1 \times 33$ \\
    \midrule
    Compact sigma proof $\Pi_{\mathrm{send}}$
      & $192$ & $192$ & $6 \times 32$ \\
    Aggregated Bulletproof
      & $754$ & $754$ & unchanged \\
    \midrule
    \textbf{Total}
      & $\mathbf{1177}$ & $\mathbf{1144}$ & $\Delta = 33$ \\
    \bottomrule
  \end{tabular}
\end{table}

\subsection{Progressive Reductions}

\Cref{tab:progressive} shows the full optimization path from the CMPT
specification baseline.
The compact sigma protocol is the primary optimisation adopted in
this document; the Bulletproofs++ upgrade~\cite{Eagen2023bppp} is an
orthogonal enhancement analysed in~\cite{Dore2026bppp}.

\begin{table}[htbp]
  \centering
  \caption{Progressive \textsf{ConfidentialMPTSend} transaction size
    reductions.  The $n = 4$ column is the production configuration;
    $n = 3$ (no auditor) is shown for comparison.
    Each step from $n = 3$ to $n = 4$ adds exactly $33$~bytes
    (one $C_{2,i}$ group element).}
  \label{tab:progressive}
  \begin{tabular}{lrrrr}
    \toprule
    Configuration & $n{=}4$ & $n{=}3$ & \multicolumn{2}{c}{vs.\ (A)} \\
    \midrule
    (A)\; Per-CT layout, separate proofs
      & $1604$ & $1604$ & \multicolumn{2}{c}{---} \\
    (B)\; Shared $C_1$, separate proofs
      & $1604$ & $1538$ & $0$/$-66$ & \\
    (B${}'$)\; (B) + merged $\Sigma$ (commit.\ form)
      & $1442$ & $1376$ & $-162$ & ($10\%$) \\
    \textbf{(B${}''$)}\; \textbf{(B) + compact $\boldsymbol{\Sigma}$}
      & $\mathbf{1177}$ & $\mathbf{1144}$ & $\mathbf{-427/-460}$ & ($\mathbf{27/29\%}$) \\
    \midrule
    (C)\; (B${}''$) + Bulletproofs++
      & $\sim 905$ & $\sim 872$ & $\sim\!-699/-732$ & ($44/46\%$) \\
    \bottomrule
  \end{tabular}
\end{table}

\noindent
Note that the $n = 4$ baseline (A) coincidentally also equals
$1604$~bytes: Cenk's shared-$C_1$ formulation with $4$ recipients
happens to match Dore's per-CT formulation with $3$ recipients
(the extra $C_{2,4}$ adds $33$~B while deduplicating $C_1$ saves
$33$~B, and the separate same-plaintext proof adds $33$~B while
losing one redundant $C_{1,i}$ saves $33$~B).

The compact sigma optimisation alone recovers $427$--$460$~bytes---more
than $60\%$ of the total potential saving including BP++.
It requires no change to the encryption scheme, no new cryptographic
primitives, and no specification of new ciphertext formats;
only the proof serialization changes.

%--------------------------------------------------------------------
\section{Extension to Other Transaction Types}
\label{sec:other-tx}
%--------------------------------------------------------------------

The preceding analysis targets \textsf{ConfidentialMPTSend}, which
carries the heaviest proof bundle.
The remaining confidential transaction types---\textsf{Convert},
\textsf{ConvertBack}, and \textsf{Clawback}---operate at the
public/confidential boundary: the transfer amount~$m$ is publicly
known, eliminating the need for same-plaintext equality proofs and,
in most cases, range proofs.
We show that compact-form techniques apply to these lighter
transactions as well.

\subsection{Convert}
\label{sec:convert}

In \textsf{ConfidentialMPTConvert}, a holder moves a publicly known
amount~$m$ from cleartext into a confidential balance.
Because $m$ appears in plaintext on-chain and the holder generates
the encryption randomness~$r$ themselves, the production CMPT
specification uses \emph{deterministic verification} rather than a
zero-knowledge proof: the transaction includes a mandatory
$32$-byte \texttt{BlindingFactor} field carrying~$r$, and
validators check
\[
  C_1 \stackrel{?}{=} r \cdot G
  \qquad\text{and}\qquad
  C_{2,i} \stackrel{?}{=} m \cdot G + r \cdot \mathit{pk}_i
\]
for each recipient key~$\mathit{pk}_i$.
No sigma proof is required or defined for Convert in the
production protocol.

For reference, the natural sigma proof for this relation---a
single Schnorr proof of knowledge of~$r$
($\Pi_{\mathrm{plain}}$)---is serialised in commitment form as
$(T_1, T_2, s) \in \G^2 \times \Zq$ costing
$2 \times 33 + 32 = 98$~bytes, or in compact form as
\[
  \Pi_{\mathrm{plain}} = (e,\, z_r) \;\in\; \Zq^2,
  \qquad\text{costing } 2 \times 32 = \mathbf{64}\text{~bytes},
\]
with reconstruction:
\begin{align}
  T_1 &= z_r \cdot G - e \cdot C_1, \label{eq:recon-conv-t1} \\
  T_2 &= z_r \cdot P_A - e \cdot (C_2 - m \cdot G).
    \label{eq:recon-conv-t2}
\end{align}
This proof would apply in contexts where the prover cannot or
should not disclose~$r$ on-chain; it is not part of the
production CMPT transaction protocol.
Clawback does not admit the r-disclosure shortcut, since the
issuer proves knowledge of~$\sk_{\mathrm{iss}}$, not randomness
(see \cref{sec:clawback}).

\subsection{Clawback}
\label{sec:clawback}

In \textsf{ConfidentialMPTClawback}, the issuer reclaims a publicly
known amount~$m$ from a holder's confidential balance.
Unlike Convert, the issuer proves knowledge of the issuer secret
key~$\sk_{\mathrm{iss}}$ rather than randomness.
Concretely, the issuer proves:
\[
  P_{\mathrm{iss}} = \sk_{\mathrm{iss}} \cdot G
  \quad\text{and}\quad
  C_2 - m \cdot G = \sk_{\mathrm{iss}} \cdot C_1,
\]
i.e.\ a Chaum--Pedersen equality-of-discrete-logs proof.
In the current commitment form this is serialised as
$(T_1, T_2, s) \in \G^2 \times \Zq$ costing $98$~bytes.

In compact form, the proof becomes
\[
  \Pi_{\mathrm{claw}} = (e,\, z_{\sk}) \;\in\; \Zq^2,
  \qquad\text{costing } 2 \times 32 = \mathbf{64}\text{~bytes},
\]
with reconstruction:
\begin{align}
  T_1 &= z_{\sk} \cdot G - e \cdot P_{\mathrm{iss}},
    \label{eq:recon-claw-t1} \\
  T_2 &= z_{\sk} \cdot C_1 - e \cdot (C_2 - m \cdot G).
    \label{eq:recon-claw-t2}
\end{align}
The Fiat--Shamir hash binds the domain tag
\texttt{CMPT\_CLAWBACK\_SIGMA}, the statement points
$P_{\mathrm{iss}}$, $C_1$, $C_2$, $m \cdot G$ (serialised as a
$33$-byte compressed point), the commitments $T_1$, $T_2$, and
the transaction context~$\TransCtx$.

\subsection{ConvertBack}
\label{sec:convertback}

In \textsf{ConfidentialMPTConvertBack}, a holder withdraws a publicly
known amount~$m$ from a confidential balance back to cleartext.
Because $m$ is public, the withdrawal ciphertext randomness~$r_w$ is
disclosed via the \texttt{BlindingFactor} field and verified
deterministically by the validator (checking
$C_{1,w} = r_w \cdot G$ and
$C_{2,w} = m \cdot G + r_w \cdot P_A$), eliminating the need for a
zero-knowledge proof of ciphertext correctness.
The sigma proof therefore covers only the balance linkage relation.
A single-value Bulletproof ($688$~bytes) then proves
$b - m \in [0, 2^{64})$.

\subsubsection*{Balance linkage relation}

The zero-knowledge proof establishes that the Pedersen commitment
$\mathit{PC}_b$ encodes the same balance~$b$ as the sender's
on-ledger balance ciphertext $(B_1, B_2)$:
\begin{equation}
  \label{eq:convertback-relation}
  \mathcal{R}_{\mathrm{bal}} = \left\{
  \bigl(\mathbf{x},\; w\bigr)
  \;\middle|\;
  \begin{aligned}
    P_A &= \sk_A \cdot G, \\
    B_2 - b \cdot G &= \sk_A \cdot B_1, \\
    \mathit{PC}_b &= b \cdot G + \rho \cdot H
  \end{aligned}
  \right\},
\end{equation}
where the public statement is
$\mathbf{x} = (P_A, B_1, B_2, \mathit{PC}_b, H)$ and the witness is
$w = (b,\, \sk_A,\, \rho) \in \Zq^3$.
This is precisely the balance sub-relation of
$\mathcal{R}_{\mathrm{send}}$, now standing alone since the amount
witness $(m, r)$ is not hidden.

\subsubsection*{AND-composed compact-form protocol}

The prover samples nonces $(t_b,\, t_{\sk},\, t_\rho)
\xleftarrow{\$} \Zq^3$ and computes three commitments:
\begin{align*}
  T_{\sk,1} &= t_{\sk} \cdot G, \\
  T_{\sk,2} &= t_b \cdot G + t_{\sk} \cdot B_1, \\
  T_b &= t_b \cdot G + t_\rho \cdot H.
\end{align*}
Challenge and responses follow the same pattern as
\cref{sec:protocol}: a single Fiat--Shamir challenge~$e$ binds all
three commitments, and responses are
$z_b = t_b + e \cdot b$,
$z_{\sk} = t_{\sk} + e \cdot \sk_A$,
$z_\rho = t_\rho + e \cdot \rho$.

The Fiat--Shamir hash is:
\begin{multline}
  \label{eq:cb-fs-hash}
  e = \mathcal{H}\!\bigl(
    \texttt{"CMPT\_CONVERTBACK\_SIGMA"} \;\|\;
    P_A \;\|\; B_1 \| B_2 \;\|\;
    \mathit{PC}_b \;\|\; \\
    T_{\sk,1} \| T_{\sk,2} \| T_b \;\|\;
    \TransCtx
  \bigr).
\end{multline}

\begin{definition}[Compact-form proof $\Pi_{\mathrm{cb}}$]
  \label{def:compact-cb}
  The non-interactive proof in compact form is
  \[
    \Pi_{\mathrm{cb}} = (e,\, z_b,\, z_\rho,\, z_{\sk})
    \;\in\; \Zq^4.
  \]
  \textbf{Verification:}
  \begin{enumerate}
    \item Reconstruct commitments:
    \begin{align}
      T_{\sk,1} &= z_{\sk} \cdot G - e \cdot P_A,
        \label{eq:recon-cb-tsk1} \\
      T_{\sk,2} &= z_b \cdot G + z_{\sk} \cdot B_1 - e \cdot B_2,
        \label{eq:recon-cb-tsk2} \\
      T_b &= z_b \cdot G + z_\rho \cdot H - e \cdot \mathit{PC}_b.
        \label{eq:recon-cb-tb}
    \end{align}
    \item Recompute $e'$ from the Fiat--Shamir hash
          (\cref{eq:cb-fs-hash}) with the reconstructed commitments.
    \item Accept iff $e' = e$.
  \end{enumerate}
\end{definition}

\noindent\textbf{Proof size.}
$4 \times 32 = \mathbf{128}$~bytes (one challenge, three responses).
The current commitment-form cost for balance linkage alone is
$195$~bytes, yielding a saving of $67$~bytes.

\begin{remark}[Security]
  The security argument for $\Pi_{\mathrm{cb}}$ is analogous to
  \cref{sec:security}: the witnesses $(b, \sk_A, \rho)$ are
  independent, so $2$-special soundness extracts each from two
  accepting transcripts with distinct challenges.
  HVZK follows from the uniform distribution of the nonces.
\end{remark}

\subsection{Cross-Transaction Summary}
\label{sec:cross-tx-summary}

\Cref{tab:cross-tx} summarises the compact-form savings across all
four confidential transaction types.

\begin{table}[htbp]
  \centering
  \caption{Sigma-proof sizes across transaction types.
    Range proofs (Bulletproof) are unchanged and shown separately.
    In production, Convert uses r-disclosure (\cref{sec:convert})
    rather than a sigma proof; the Convert row shows hypothetical
    compact-form sizes for reference.}
  \label{tab:cross-tx}
  \begin{tabular}{l@{\quad}r@{\quad}r@{\quad}r@{\quad}r}
    \toprule
    Transaction type
    & \shortstack{Sigma\\(current)}
    & \shortstack{Sigma\\(compact)}
    & \shortstack{Range\\proof}
    & \shortstack{Sigma\\saving} \\
    \midrule
    \textsf{Send} ($n{=}4$)
      & $619$ & $192$ & $754$ & $427$ ($69\%$) \\
    \textsf{ConvertBack}
      & $195$ & $128$ & $688$ & $67$ ($34\%$) \\
    \textsf{Convert}
      & $98$ & $64$ & --- & $34$ ($35\%$) \\
    \textsf{Clawback}
      & $98$ & $64$ & --- & $34$ ($35\%$) \\
    \bottomrule
  \end{tabular}
\end{table}

\noindent
The absolute savings are largest for \textsf{Send} ($427$~bytes), as
expected given its richer proof structure.
For \textsf{ConvertBack}, the $67$-byte reduction yields a total proof
payload of $128 + 688 = 816$~bytes, down from
$195 + 688 = 883$~bytes ($8\%$).
Convert and Clawback see modest absolute savings ($34$~bytes each),
though the percentage reduction in sigma overhead is comparable.

The compact-form technique is uniform: the same design
pattern---AND-compose all sigma relations, transmit $(e, \mathbf{z})$
instead of $(\mathbf{T}, \mathbf{z})$, reconstruct commitments during
verification---applies to every transaction type without exception.

\subsection{Implementation Changes}
\label{sec:other-tx-impl}

Adopting compact sigma for the non-Send transaction types requires
coordinated changes in the cryptographic library
(\texttt{mpt-crypto}) and the ledger node (\texttt{rippled}).

\subsubsection*{Changes in \texttt{mpt-crypto}}

The \texttt{mpt-crypto} library already implements each building-block
sigma protocol as a separate prove/verify pair:
\texttt{secp256k1\_equality\_plaintext\_\{prove,verify\}} for the
plaintext~PoK ($\Pi_{\mathrm{plain}}$, $98$~bytes) and
\texttt{secp256k1\_elgamal\_pedersen\_link\_\{prove,verify\}} for
the linkage proof ($\Pi_{\mathrm{link}}$, $195$~bytes).
The different transaction types compose these building blocks:
Convert uses deterministic r-disclosure (no sigma proof; see
\cref{sec:convert}); Clawback uses a Chaum--Pedersen
equality-of-discrete-logs proof for~$\sk_{\mathrm{iss}}$;
ConvertBack uses
$\Pi_{\mathrm{link\text{-}B}}$ plus a
Bulletproof; Send uses all three plus the shared-$r$ equality proof.

The compact sigma implementation for Send (PR~\#18) establishes
the pattern: a single new source file
(\texttt{proof\_compact\_standard.c}) with dedicated
prove/verify entry points, a size constant, and a
Fiat--Shamir domain tag distinct from the commitment-form proofs.
No changes to the underlying scalar or group-arithmetic layer are
needed.

Convert requires no new source file: the production protocol uses
deterministic r-disclosure via \texttt{BlindingFactor}, so no sigma
prove/verify pair is defined.
Following this convention, we propose one new source file per
remaining transaction type:

\begin{description}[style=nextline,leftmargin=1.5em]
  \item[Clawback (compact $\Pi_{\mathrm{claw}}$).]
    Implemented in \texttt{proof\_compact\_clawback.c}.
    The proof shrinks from $98$ to $64$~bytes.
    The prover samples one nonce~$t_{\sk}$, computes
    $z_{\sk} = t_{\sk} + e \cdot \sk_{\mathrm{iss}}$, and serialises
    $(e, z_{\sk})$.
    The verifier reconstructs $T_1, T_2$ via
    \crefrange{eq:recon-claw-t1}{eq:recon-claw-t2} and re-hashes
    under the domain tag \texttt{CMPT\_CLAWBACK\_SIGMA}.
    Note that $m \cdot G$ enters the hash as a $33$-byte compressed
    point, not a raw scalar.

  \item[ConvertBack (compact $\Pi_{\mathrm{cb}}$).]
    A new source file (e.g.\
    \texttt{proof\_compact\_convertback.c}) providing
    prove/verify entry points.
    Proof size: $128$~bytes.
    Because the withdrawal ciphertext randomness~$r_w$ is
    disclosed via the \texttt{BlindingFactor} field and verified
    deterministically, the sigma proof covers only the balance
    linkage relation (\cref{sec:convertback}).
    The implementation samples three nonces, computes three
    commitments, derives~$e$ under the domain tag
    \texttt{CMPT\_CONVERTBACK\_SIGMA}, computes three responses,
    and serialises $(e, z_b, z_\rho, z_{\sk})$.

    The Bulletproof range proof for
    $b - m \in [0, 2^{64})$ is unchanged;
    the existing Bulletproof prove/verify functions continue to
    operate on the Pedersen commitment
    $\mathit{PC}_b - m \cdot G$.
\end{description}

\noindent
PR~\#18 also introduces \texttt{mpt\_internal.h} with shared
helpers (\texttt{generate\_random\_scalar},
\texttt{mpt\_uint64\_to\_scalar}, etc.) that the new source
files can reuse directly.
In all cases, the old commitment-form functions on main are
retained; the compact variants are strictly additive.
The nonce-generation guidance of \cref{sec:nonces} (deterministic
synthetic nonces incorporating the full witness) applies equally to
ConvertBack and Clawback.

\subsubsection*{Changes in \texttt{rippled}}

The ledger node calls \texttt{mpt-crypto} during transaction
validation.
Required changes:

\begin{enumerate}
  \item \textbf{Proof-size constants.}
    The serialisation layer must accept the new proof sizes:
    $64$~bytes for Convert and Clawback sigma proofs and $128$~bytes for
    ConvertBack sigma proofs (replacing $98$ and $195$
    respectively).
    These constants gate deserialisation and buffer allocation.

  \item \textbf{Verification call sites.}
    Each transaction handler currently calls the commitment-form
    verify function.
    For \textsf{ConfidentialMPTConvertBack}, the two separate verify
    calls (plaintext proof and linkage proof) are replaced by a
    single call to the compact ConvertBack verifier.
    For Convert, a single call replaces the existing one with the
    only visible change being the proof buffer size.
    For Clawback, the call similarly replaces the existing one, but
    the underlying proof relation changes from knowledge of
    randomness to knowledge of~$\sk_{\mathrm{iss}}$.

  \item \textbf{TransactionContextID construction.}
    The Fiat--Shamir hash in each compact proof binds the
    $\TransCtx$.
    The \texttt{rippled} code that constructs these context IDs is
    transaction-type-specific and is already in place; no change is
    needed unless the domain tag is incorporated into the
    context ID itself rather than the proof hash (the current
    design keeps them separate).

  \item \textbf{Amendment gating.}
    Because the proof format is a consensus-visible change (validators
    must agree on the serialisation), the new compact formats must be
    gated behind an XRPL amendment.
    The amendment can cover all three transaction types simultaneously,
    or be staged (Send first, others in a follow-up), depending on
    rollout preferences.
    During the transition, validators must accept \emph{both}
    commitment-form and compact-form proofs until the amendment
    activates, at which point only compact form is accepted.
\end{enumerate}

\begin{remark}[Scope relative to Send]
  The compact sigma for \textsf{Send} is the higher-priority
  implementation target, since it delivers the largest absolute saving
  ($427$~bytes) on the most frequent transaction type.
  The Convert/ConvertBack/Clawback extensions described here are
  incremental: they reuse the same architectural pattern and can share
  infrastructure (e.g., the scalar serialisation helpers in
  \texttt{mpt\_internal.h}) already introduced for Send.
  They can be implemented in a follow-up PR once the Send compact
  sigma has been validated.
\end{remark}

%--------------------------------------------------------------------
\appendix
\section{Twisted ElGamal Variant}
\label{app:twisted}
%--------------------------------------------------------------------

We briefly present the compact sigma protocol for the twisted ElGamal
construction of~\cite{Dore2026twisted} and analyse why its advantage
over standard EC-ElGamal shrinks under compact form.
The team has elected \textbf{not} to adopt twisted ElGamal for the
CMPT specification, for the reasons summarised in
\cref{sec:twisted-assessment}.

\subsection{Combined Language}

In twisted ElGamal, each ciphertext component $Y = r \cdot G + m \cdot H$
simultaneously serves as a Pedersen commitment, eliminating separate
$\mathit{PC}_m$ and $\mathit{PC}_b$.
The refreshed remainder ciphertext $(Y^*, X^*)$ uses a DDH proof
to link to the balance.
The combined relation is:
\begin{multline*}
  L_{\mathrm{comb}}^{(n)} = \bigl\{
  (\pk_i, Y, X_i, \pk_A, D_X, D_Y) \;\big|\;
  \exists\,(r, m, \sk_A) \in \Zq^3: \\
  Y = r \cdot G + m \cdot H
  \;\wedge\;
  X_i = r \cdot \pk_i\; \forall i
  \;\wedge\;
  \pk_A = \sk_A \cdot G
  \;\wedge\;
  D_X = \sk_A \cdot D_Y
  \bigr\},
\end{multline*}
where $D_X, D_Y$ are the difference ciphertext components.
The witness is $(r, m, \sk_A)$---only three scalars, compared to five
for standard EC-ElGamal.

\subsection{Compact-Form Protocol}

The AND-composed sigma protocol has nonces $(a, b, k)$ for
$(r, m, \sk_A)$, with $n + 3$ first-round commitments:
\begin{align*}
  T_Y &= a \cdot G + b \cdot H, &
  T_{X,i} &= a \cdot \pk_i \quad (i = 1, \ldots, n), \\
  K_1 &= k \cdot G, &
  K_2 &= k \cdot D_Y.
\end{align*}
Responses: $z_1 = a + e \cdot r$, $z_2 = b + e \cdot m$,
$z_3 = k + e \cdot \sk_A$.

\medskip
\noindent\textbf{Compact form.}
$\Pi_{\mathrm{comb}} = (e, z_1, z_2, z_3) \in \Zq^4$:
$4 \times 32 = \mathbf{128}$~bytes.

\medskip
\noindent\textbf{Verification.}
Reconstruct commitments:
\begin{align*}
  T_Y &= z_1 \cdot G + z_2 \cdot H - e \cdot Y, &
  T_{X,i} &= z_1 \cdot \pk_i - e \cdot X_i, \\
  K_1 &= z_3 \cdot G - e \cdot \pk_A, &
  K_2 &= z_3 \cdot D_Y - e \cdot D_X.
\end{align*}
Check $e \stackrel{?}{=} \mathcal{H}(\cdots)$.

Security follows from the same arguments as \cref{sec:security},
with the witness and nonce spaces having dimension~$3$ instead of~$5$.

\subsection{The 64-Byte Gap}
\label{sec:twisted-gap}

\Cref{tab:twisted-comparison} compares the sigma overhead of both
schemes across proof formats.

\begin{table}[htbp]
  \centering
  \caption{%
    Sigma-proof size comparison: standard vs.\ twisted EC-ElGamal.
    Commitment-form sizes depend on~$n$; compact-form sizes do not.}
  \label{tab:twisted-comparison}
  \begin{tabular}{l@{\quad}rr@{\;\;}rr@{\;\;}r}
    \toprule
    & \multicolumn{2}{c}{\shortstack{Separate\\(commit.\ form)}}
    & \multicolumn{2}{c}{\shortstack{Merged\\(commit.\ form)}}
    & \shortstack{Compact\\form} \\
    & $n{=}4$ & $n{=}3$
    & $n{=}4$ & $n{=}3$
    & \\
    \midrule
    \textbf{Std.\ EC-ElGamal}
    & $619$ & $586$
    & $457$ & $424$
    & $192$ \\[2pt]
    \textbf{Twisted ElGamal}
    & $327$ & $294$
    & $327$ & $294$
    & $128$ \\
    \midrule
    \textbf{Std.\ $\to$ Twisted savings}
    & $292$ & $292$
    & $130$ & $130$
    & $\mathbf{64}$ \\
    \bottomrule
  \end{tabular}
\end{table}

\noindent
Two observations:

\begin{enumerate}
  \item \textbf{Compact form benefits standard EC-ElGamal more.}
    Standard EC-ElGamal has more group-element commitments to
    eliminate ($8$ vs.\ $6$ in merged form), yielding a larger
    absolute reduction ($-232$~B vs.\ $-166$~B).

  \item \textbf{The twisted ElGamal advantage shrinks from $292$ to
    $64$~bytes.}
    In commitment form, twisted ElGamal saves bytes primarily by
    eliminating linkage proofs whose cost comes from group-element
    commitments ($33$~B each).
    Compact form removes \emph{all} group-element commitments from
    \emph{both} schemes.
    What remains is the irreducible per-witness cost: each independent
    scalar witness contributes one $32$-byte response.
    Twisted ElGamal has three witnesses $(r, m, \sk_A)$ whereas
    standard EC-ElGamal has five $(m, r, b, \rho, \sk_A)$, because
    the dual-purpose ciphertext/commitment structure of twisted
    ElGamal makes separate Pedersen commitments---and their associated
    witnesses $\rho$ and~$b$---unnecessary.
\end{enumerate}

\noindent
The $64$-byte gap is \textbf{structural and irreducible}: it comes
from the number of independent scalar witnesses, not from group
elements.
No further compression technique---including the $\mathcal{O}(1)$
equality proof of \cref{sec:equality-o1}---can close it.
Adding or removing auditors changes the ciphertext count by
$\pm 33$~bytes per recipient, identically for both schemes, so the
relative comparison is unaffected by~$n$.

\subsection{Full Transaction Comparison}

\Cref{tab:twisted-full} shows the complete transaction sizes.

\begin{table}[htbp]
  \centering
  \caption{%
    Full \textsf{ConfidentialMPTSend} comparison under compact sigma.
    The $\Delta$ column (twisted vs.\ standard) is $n$-independent.}
  \label{tab:twisted-full}
  \begin{tabular}{l@{\quad}rr@{\quad}rr@{\quad}r}
    \toprule
    & \multicolumn{2}{c}{Std.\ EG compact}
    & \multicolumn{2}{c}{Twisted EG compact}
    & \\
    Component
      & $n{=}4$ & $n{=}3$
      & $n{=}4$ & $n{=}3$
      & $\Delta$ \\
    \midrule
    Ciphertext / commitments
      & $231$ & $198$
      & $231$ & $198$
      & $0$ \\
    Compact sigma proof
      & $192$ & $192$
      & $128$ & $128$
      & $-64$ \\
    Aggregated Bulletproof
      & $754$ & $754$
      & $754$ & $754$
      & $0$ \\
    \midrule
    \textbf{Total}
      & $\mathbf{1177}$ & $\mathbf{1144}$
      & $\mathbf{1113}$ & $\mathbf{1080}$
      & $\mathbf{-64}$ \\
    \midrule
    \textit{With BP++ upgrade}
      & $\sim 905$ & $\sim 872$
      & $\sim 841$ & $\sim 816$
      & $\sim\!-64$ \\
    \bottomrule
  \end{tabular}
\end{table}

\subsection{Assessment}
\label{sec:twisted-assessment}

The $64$-byte advantage of twisted ElGamal must be weighed against its
costs:

\begin{enumerate}[label=(\roman*)]
  \item \textbf{New API surface.}
        Twisted ElGamal requires a generator-swap API
        ($\mathit{value\_gen} = H$, $\mathit{blinding\_gen} = G$)
        for both Pedersen commitments and Bulletproof range proofs.
        Passing generators in the wrong order silently produces
        insecure proofs.

  \item \textbf{Audit scope.}
        The existing \texttt{mpt-crypto} library implements standard
        EC-ElGamal with substantial test coverage and is currently
        under audit.
        Twisted ElGamal requires auditing new encryption, decryption,
        and proof primitives.

  \item \textbf{Specification changes.}
        The CMPT spec and XLS-96 are written for standard EC-ElGamal.
        Switching requires updating both documents, re-specifying all
        Fiat--Shamir transcripts, and revising the on-chain ciphertext
        format.
\end{enumerate}

\noindent
For context, the BP++ upgrade saves $\sim 272$~bytes---more than
$4\times$ the twisted ElGamal advantage.
At $n = 4$, standard EC-ElGamal with compact sigma and BP++ gives
$\sim 905$~bytes, compared to $\sim 841$~bytes for twisted ElGamal
with the same optimizations: a difference of $64$~bytes ($7\%$) on a
total already below $1$~KB.

\begin{remark}[Recommendation]
  Given that the compact sigma optimization is adopted, the
  $64$-byte saving from twisted ElGamal does not justify the
  specification, implementation, and audit costs.
  The team's decision to proceed with standard EC-ElGamal and
  compact sigma is sound.
\end{remark}

%--------------------------------------------------------------------
\bibliographystyle{alpha}

\providecommand{\etalchar}[1]{$^{#1}$}
\begin{thebibliography}{EKRN23}

\bibitem[BN06]{BellareNeven2006}
Mihir Bellare and Gregory Neven.
\newblock Multi-signatures in the plain public-key model and a general
  forking lemma.
\newblock In {\em ACM CCS 2006}, pages 390--399, 2006.

\bibitem[BBB{\etalchar{+}}18]{Bunz2018bulletproofs}
Benedikt B{\"u}nz, Jonathan Bootle, Dan Boneh, Andrew Poelstra, Pieter Wuille,
  and Greg Maxwell.
\newblock Bulletproofs: Short proofs for confidential transactions and more.
\newblock In {\em IEEE Symposium on Security and Privacy}, pages 315--334,
  2018.

\bibitem[CDS94]{CramerDamgardSchoenmakers1994}
Ronald Cramer, Ivan Damg{\aa}rd, and Berry Schoenmakers.
\newblock Proofs of partial knowledge and simplified design of witness hiding
  protocols.
\newblock In {\em CRYPTO~'94}, volume 839 of {\em LNCS}, pages 174--187.
  Springer, 1994.

\bibitem[CMA26]{CenkMalhotraAkinyele2026}
Murat Cenk, Aanchal Malhotra, and Ayo Akinyele.
\newblock Confidential multi-purpose tokens on the {XRP} ledger.
\newblock Ripple Research, February 2026.
\newblock Internal specification, version 2026-02-01.

\bibitem[Cen26]{Cenk2026unification}
Murat Cenk.
\newblock Proof unification and transcript compression for efficient
  confidential transfers.
\newblock Ripple Research, March 2026.
\newblock Internal working document.

\bibitem[Dor26a]{Dore2026twisted}
Tess Dore.
\newblock Twisted {ElGamal} and shared randomness for confidential
  multi-purpose tokens on {XRPL}.
\newblock RippleX Research, March 2026.

\bibitem[Dor26b]{Dore2026bppp}
Tess Dore.
\newblock Bulletproofs++ upgrade for confidential multi-purpose tokens
  on {XRPL}.
\newblock RippleX Research, March 2026.

\bibitem[EKRN23]{Eagen2023bppp}
Liam Eagen, Sanket Kanjalkar, Tim Ruffing, and Jonas Nick.
\newblock Bulletproofs++: Next generation confidential transactions via
  reciprocal set membership arguments.
\newblock Cryptology ePrint Archive, Report 2022/510, 2023.
\newblock Revised July 2023. \url{https://eprint.iacr.org/2022/510}.

\bibitem[Por13]{RFC6979}
Thomas Pornin.
\newblock Deterministic usage of the {Digital Signature Algorithm (DSA)}
  and {Elliptic Curve Digital Signature Algorithm (ECDSA)}.
\newblock {RFC}~6979, August 2013.
\newblock \url{https://www.rfc-editor.org/rfc/rfc6979}.

\bibitem[Sch91]{Schnorr1991}
Claus-Peter Schnorr.
\newblock Efficient signature generation by smart cards.
\newblock {\em Journal of Cryptology}, 4(3):161--174, 1991.

\end{thebibliography}

\end{document}
